\chapter{Supplementary Results and Sensitivity Analyses}
\label{app:supplementary-results}

\section{Matched Two-Epoch Representation Diagnosis}
The preceding four-arm study compared sixteen and thirty-two frames under unequal epoch
budgets. Table~\ref{tab:fourarm-visual-supplement} retains its results separately from
the later matched diagnostic. All four arms used the same 4,530 actions, five outer
folds and seed family~42.

\begin{table}[htbp]
\caption{Earlier Four-Arm Visual Comparison on the Common REASSEMBLE Cohort}
\label{tab:fourarm-visual-supplement}
\centering\small
\begin{tabular}{lrrrrr}\toprule
CNN state & Frames & Epochs & Visual AP & Sensor AP & F2 AP\\\midrule
Frozen & 16 & 5 & 0.221644 & 0.420803 & 0.440680\\
Final-stage tuned & 16 & 2 & 0.224466 & 0.420803 & 0.437521\\
Frozen & 32 & 5 & 0.219427 & 0.420803 & 0.439584\\
Final-stage tuned & 32 & 2 & 0.230609 & 0.420803 & 0.438566\\\bottomrule
\end{tabular}
\displaynote{Failure-positive pooled AP. Frozen arms fitted the head; tuned arms jointly
fitted the head and final CNN stage. Each arm completed twenty-five visual fits. Unequal
epoch budgets do not isolate the effect of adaptation. The common follow-up sensor
reference differs from the core selected-$C$ baseline (AP 0.4294).}
\end{table}
\ev{Original four-arm results and fitting scope}{\EvidenceRef{E144}{},
\EvidenceRef{E145}{}}

The frozen32 versus fine-tuned32 diagnosis then compares equal epoch and data budgets on
the same complete 4,530-action cohort. It contains five matched outer folds for seed~42
and five for seed~20261023, with 75 new and 25 reused visual fits. All reused fits belong
to the earlier seed-42 two-epoch fine-tuned32 arm: four inner fits and one final fit per
outer fold. They are not an independent replication. Equal budgets do not imply equal
parameter counts, computation or training time. The fixed numerical branch is shared
between arms, so a visual or fusion difference does not arise from refitting a different
sensor model.\ev{Reused fit identities}{\EvidenceRef{E143}{}}

Table~\ref{tab:twoepoch-supplement} gives pooled failure-positive AP for each method at
both seeds. The intervals resample paired recording clusters with the models held fixed.
Mean-fold AP is a separate summary that weights folds differently.

\begin{table}[htbp]
\caption{Fine-Tuned Minus Frozen Average Precision After Two Epochs}\label{tab:twoepoch-supplement}
\centering\small

\begin{tabular}{llrrl}\toprule
Seed & Method & Frozen & Fine-tuned & Difference [95\% interval]\\\midrule
42 & Visual &0.222708&0.230609&+0.007902 [$-0.017312$, $+0.033219$]\\
42 & Averaging &0.444720&0.440946&$-0.003774$ [$-0.014238$, $+0.007851$]\\
42 & F2 &0.439363&0.438566&$-0.000796$ [$-0.006424$, $+0.006072$]\\
20261023 & Visual &0.216799&0.200517&$-0.016282$ [$-0.039636$, $+0.009790$]\\
20261023 & Averaging &0.436495&0.428489&$-0.008006$ [$-0.018180$, $+0.003279$]\\
20261023 & F2 &0.435202&0.432380&$-0.002822$ [$-0.008285$, $+0.003602$]\\\bottomrule
\end{tabular}
\displaynote{The conditional 95\% intervals quantify contrasts between the fitted
predictors. Variation across training seeds is a separate source of uncertainty.}
\end{table}
\ev{Matched two-epoch results}{\EvidenceRef{E060}{}}

At seed~42, the small gain in visual AP coincided with lower AP for both fusion methods.
At the second seed, fine-tuning reduced the point estimate for all three methods. Every
interval includes zero, so the direction of an adaptation effect remains uncertain.

The diagnostic compared training duration, complementary errors, branch calibration, F2
fitting and the difference between out-of-fold and assessment scores. The available
comparisons did not identify a cause for the variation across seeds. Class-conditional
Brier score and log loss characterise prediction errors on actions the sensor
misclassified, but cannot determine whether those images contained sufficient visual
evidence.\ev{Unresolved causal diagnosis}{\EvidenceRef{E060}{}}

\section{The Same-Test 753-Action Follow-Up}
The 753-action cohort is exactly original outer fold~0, not an independent external
dataset. It contains 77 failures from 26 recordings within the full 4,530-action,
148-recording cohort. The correspondence is established through native recording/segment
identities, timestamps and labels, with older sample hashes retained as aliases. Failure
prevalence is therefore different from that of the pooled five-fold assessment.
C20 denotes the twenty-epoch full-CNN condition.

\begin{table}[htbp]
\caption{Prediction Performance on the Same 753 Held-Out Actions}\label{tab:753-supplement}
\centering\small

\begin{tabular}{lr}\toprule Model/recipe&AP\\\midrule
Original frozen visual&0.232805\\
C20 full-CNN outer refit, seed42&0.215321\\
Frozen CNN100&0.193776\\
Full CNN100&0.212585\\
Static sensor, original score vector&0.327705\\
Original F2 restricted to fold0&0.367155\\
Frozen CNN100 plus F2&0.334977\\
Full CNN100 plus F2&0.349094\\
Frozen CNN100 plus averaging&0.348007\\
Full CNN100 plus averaging&0.358533\\\bottomrule
\end{tabular}
\displaynote{Models were evaluated on identical actions with intact inputs. The
training, calibration, and stacking procedures differed between generations.}
\end{table}
\ev{Matched cohort reconciliation}{\EvidenceRef{E050}{}}

Fine-tuning improves the hundred-epoch visual recipe relative to its frozen
hundred-epoch control, but the original frozen recipe has a higher point AP on the same
actions. Similarly, full100 F2 exceeds frozen100 F2, while the existing paired interval
crosses zero. The original F2 row uses a different nested stacking protocol. Neither
comparison establishes that the evolving later recipe is the best general architecture.

The selected inner job originally fitted 1,752 training actions, selected on 419
validation actions and calibrated on 686. Its original assessment contained 920 inner
actions. The hundred-epoch follow-up retains those fitting roles but assesses the 753
original outer actions; the 920-action bank becomes the held-out fusion-fitting bank. In
contrast, the C20 outer refit uses 2,385/581/811 training/validation/calibration
actions. These role differences remain even after the assessment identities match.

\section{Aggregation, Aliases and Sensitivity Boundaries}
Across all five original outer folds, pooled AP was 0.230658 for vision, 0.429392 for
the static sensor and 0.447299 for F2. The C20 full-CNN model achieved pooled AP of
0.312127. The hundred-epoch follow-up assessed only 753 of those actions and produced no
predictions for the other 3,777.

Several apparent comparisons reuse identical predictions. The static-sensor comparison reuses the
original sensor predictions, and the hundred-epoch fusion study uses the standalone visual
outputs. Native recording and segment keys assign the 753 outer assessment actions to 26
recordings. The 28-recording list in the selected-fit metadata refers to the inner
prediction bank.

MSM varied the target definition and grouping policy and compared exposure to synthetic
corruption. HANDOVER evaluated its baselines without fitting learned fusion because
participant-disjoint out-of-fold predictions were unavailable. ViFailback compared the
original frozen representation and later matched donor conditions. Differences in
targets, prevalence and grouping limit any comparison of average precision across these
datasets.\ev{Sensitivity and status inventory}{\EvidenceRef{E106}{}}

\section{Saved-Prediction Check of the Core F2 Visual Term}
The retained saved-prediction check verified 4,530 unique actions, their 509 failure
labels, fold assignments and separation of the 148 recording groups. The incremental-fusion experiment
and the subsequent fusion-method comparison used identical held-out branch scores;
the sensor scores also matched the separately saved statistical control. The original
unimodal visual prediction file was unavailable for this check, so its identity could
not be rechecked directly. The saved-prediction calculation did not reproduce training.
The fitting implementation and metadata place the visual column before the sensor column.

All visual coefficients are positive (3.099--4.063), as are all sensor coefficients
(6.218--6.709). Reconstructing F2 from them differs from the saved probabilities by at
most $2.22\times10^{-16}$. Independently calculated threshold-group AP and
positive--negative ordering changes agree with the library metrics. Omitting the visual
term preserves each fold's sensor AP and AUROC; Table~\ref{tab:f2-fold-audit} gives the
full-model changes.

\begin{table}[htbp]
\caption{Core F2 Minus Sensor Ranking Within Each Held-Out Fold}\label{tab:f2-fold-audit}
\centering\small
\begin{tabular}{rrrrrr}\toprule
Fold & Actions & Failures & Recordings & $\Delta$AP & $\Delta$AUROC\\\midrule
0 & 753 & 77 & 26 & $+0.039451$ & $+0.020691$\\
1 & 964 & 109 & 31 & $+0.014616$ & $+0.011674$\\
2 & 1,037 & 125 & 32 & $+0.011329$ & $+0.024430$\\
3 & 906 & 99 & 29 & $-0.004689$ & $+0.002891$\\
4 & 870 & 99 & 30 & $+0.036423$ & $+0.020045$\\\bottomrule
\end{tabular}
\displaynote{Failure-positive metrics on saved outer predictions. Folds share training
data and are not independent experimental replications. No new training or interval
estimation was performed.}
\end{table}
Fold~3 has 3,051 corrected and 2,820 damaged positive--negative pairs. Its positive
AUROC change and negative AP change are consistent because AP weights retrieval
positions differently. Pooled AUROC is 0.779325 for sensors, 0.777037 after arithmetic
omission and 0.794968 for F2. Fold-dependent score alignment therefore affects pooled
metrics without changing within-fold sensor ordering. Setting the visual probability to
zero is a particular constant reference, not an identified decomposition of pooled
importance. The reduced map retains jointly fitted parameters; it is not a refitted
sensor-only comparator.\ev{Saved-prediction verification recorded on 25 November 2026}{\EvidenceRef{E111}{}, \EvidenceRef{E120}{}}

\section{Historical PHM Ridge Reference and Target Boundary}
\label{sec:phm-ridge-supplement}
The retained Ridge run-summary predictor has raw top-three target MAE 5.0151 percentage
points on eight EXP-F runs in the historical EXP-B/EXP-A/EXP-F fixed split. Its
causal-monotonic target row instead gives 4.9357; those outcomes are not interchangeable.
The target-version study identifies this as the v2 comparison, whose training-mean
reference was 0.6797. The run's saved metric rows verify the reported numbers; they do
not establish a v3 LOEO Ridge result or complete training-source lineage. The main
chapter therefore keeps the Ridge recipe supplementary and omits its configuration-only
alpha grid.\ev{Saved Ridge metrics and target history}{\EvidenceRef{E138}{},
\EvidenceRef{E021}{}}

\section{PHM Four-Arm Diagnostic With Preprocessing Exposure}
\label{sec:phm-fourarm-supplement}
The separate v3 four-arm report gives three-seed mean run-level MAEs of 0.8288 for
modular scalar fusion, 0.9074 for process-conditioned gating and 0.9643 for
concatenation, against a fold-specific training-mean reference of 0.8042 percentage
points. These are retained reported diagnostics. As explained in
Section~\ref{sec:phm-diagnostics}, preprocessing was fitted before rotating the folds,
allowing held-out EXP-B measurements to influence scaling. These values therefore do
not constitute an isolated LOEO comparison and are not combined with the principal
fold-local result in Table~\ref{tab:phm-loeo}. They do not establish that the sensors
contain no predictive information.\ev{Reported four-arm values and preprocessing scope}
{\EvidenceRef{E129}{}, \EvidenceRef{E027}{}}
